%lezione 16 del 22 aprile 2026 pag 35 appunti
%-----------------------------------------------------------------------------------------------------
I coefficienti $A,C\in\mathbb{C}$, mentre i $\sin$ sono reali, dunque
$$|A|e^{i\varphi_A}\sin(\alpha)=|C|e^{i\varphi_C}\sin(\gamma) \rightarrow \varphi_A = \varphi_C
\phantom{x}\lor\phantom{x}\varphi_A = \varphi_C+\pi$$
inoltre, il modulo del LHS è un numero nell'intervallo $[-|C|,|C|]$, mentre quello del RHS in
$[-|A|,|A|]$, dunque $|A|=|C|$ e si trova
$$A,C\in\mathbb{R},\phantom{x}A=C\phantom{x}\lor\phantom{x}A=-C$$
I due casi, ricordando che $\alpha\not=\gamma$, implicano che
$$
\begin{dcases}
    A=C\\
    A=-C
\end{dcases}
\rightarrow
\begin{dcases}
\sin(\alpha)=\sin(\gamma) \\
\sin(\alpha)=-\sin(\gamma)
\end{dcases}
\rightarrow
\begin{dcases}
    \gamma=-\alpha+(2n+1)\pi&\rightarrow\alpha+\gamma=(2n+1)\pi\\ 
    \gamma=-\alpha+2n\pi&\rightarrow\alpha+\gamma = 2n\pi
\end{dcases}
$$
Si ha dunque $\alpha+\gamma=(n+1)\pi$ con $n=0,1,...$, esplicitando $\alpha$ e $\gamma$ si trova infine
\begin{gather*}
\alpha+\gamma = \int_{b}^{x}dx'k(x') + \int_{x}^{a}dx'k(x') +\frac{\pi}{2} = (n+1)\pi \rightarrow\\
\rightarrow\boxed{\int_{b}^{a}dx'k(x') = \frac{\pi}{2}+n\pi,\phantom{x}n=0,1,...}
\end{gather*}
che è una condizione di quantizzazione, come ci si aspetta da un problema legato.

Classicamente, $(b,a)$ è la metà di un periodo di oscillazione del sistema, quindi
$$\int_{b}^{a}dx'k(x') = \frac{2\pi}{2}\oint dx'\frac{p}{h} = \frac{\pi}{2}+n\pi \rightarrow
\oint pdx' = \left(n+\frac{1}{2}\right)\pi
$$
che è la condizione di quantizzazione di Sommerfield.

\textit{Osservazione} L'approssimazione WKB vale per $p$ grandi, quindi per $\lambda$ di De Broglie
piccole: in questo caso deve valere $\lambda<<b-a$, che vuol dire $n$ grande, perché esso è il numero
di oscillazioni che la $\psi$ compie nella zone classicamente permessa $(b,a)$.

In questo modo, si può anche utilizzare $\int_{b}^{a}dx'k(x')$ come discriminante tra potenziali per
cui vale l'approssimazione WKB.

\textbf{Esempio} L'oscillatore armonico

$V(x)=\frac{1}{2}m\omega^2x^2$, con energia $E=\frac{1}{2}m\omega^2a^2$ e punti di inversione classici
$\pm a$:
$$k(x) = \frac{1}{\hbar}\sqrt{2m\frac{1}{2}m\omega^2(a^2-x^2)} = \frac{m\omega}{\hbar}\sqrt{a^2-x^2}$$
la condizione di quantizzazione è
$$\int_{-a}^{a}dxk(x)=\frac{m\omega}{\hbar}dx\sqrt{a^2-x^2} = \frac{m\omega a^2}{\hbar}
\int_{-\frac{\pi}{2}}^{\frac{pi}{2}}d\theta\cos^2(\theta) = \frac{m\omega a^2\pi}{2\hbar} =
\left(n+\frac{1}{2}\right)\pi
$$
Riconoscendo $E=\frac{1}{2}m\omega^2a^2$ si ha
$$\frac{E}{\hbar\omega} = n+\frac{1}{2} \rightarrow E = \left(n+\frac{1}{2}\right)\hbar\omega$$
che è lo spettro energetico completo e non un'approssimazione perché l'oscillatore armonico è
quadratico.
